\documentclass[../../../include/open-logic-section]{subfiles}

\begin{document}

\olfileid{sth}{cardinals}{zfc}
\olsection{$\ZFC$: একটি মাইলফলক}

সুক্রমবিন্যাস যোগ করার ফলে তত্ত্বের শেষ মাইলফলকে পৌঁছেছি।
এখন $\ZFC$-র জন্য প্রয়োজনীয় সব স্বতঃসিদ্ধ আমাদের আছে।
নির্দিষ্টভাবে:

\begin{defn}
$\ZFC$ তত্ত্বের স্বতঃসিদ্ধগুলি হলো: সদস্যভিত্তিক সমতা, সংযোগ,
যুগল, ঘাত সেট, অসীমতা, ভিত্তি, সুক্রমবিন্যাস এবং পৃথকীকরণ ও
প্রতিস্থাপন ছকের সব দৃষ্টান্ত। অন্যভাবে বললে, $\ZF$-এর
সঙ্গে সুক্রমবিন্যাস যোগ করলেই $\ZFC$ পাওয়া যায়।
\end{defn}

$\ZFC$-র শেষ অক্ষরটি \emph{চয়ন}-সহ \emph{জার্মেলো--ফ্রেঙ্কেল}
সেটতত্ত্বের নির্দেশক। এটি প্রথমে একটু অদ্ভুত লাগতে পারে:
যোগ করা স্বতঃসিদ্ধের নাম তো ``সুক্রমবিন্যাস'', ``চয়ন'' নয়।
কিন্তু পরে \emph{চয়ন} সূত্রবদ্ধ করলে দেখা যাবে, $\ZF$-এর
সাপেক্ষে সুক্রমবিন্যাস ও চয়ন তুল্য
(দেখুন \olref[choice][woproblem]{thmwochoice})। তাই কোনটিকে
``মূল'' স্বতঃসিদ্ধ ধরব, তা নিয়ে ফলের কোনো পার্থক্য নেই।
আর ``$\ZFC$'' নামটি গবেষণাসাহিত্যে সর্বজনস্বীকৃত।

\end{document}
